\section{Complete Stieltjes reflection through a positive mass}
\label{qspec:sec:stieltjes}

Define the following ordinary, absolutely convergent sums:
\begin{align}
D_m^+(a,b)=\sum_{n\ge0}\biggl\{&
\frac{(-\log Q_b(x_n))^m}{\sqrt{Q_b(x_n)}}
-\frac{(-2\log x_n)^m}{x_n}\notag\\[-2pt]
&+\frac{(-\log Q_b(y_n))^m}{\sqrt{Q_b(y_n)}}
-\frac{(-2\log y_n)^m}{y_n}\biggr\},\label{qspec:eq:Dplus}\\
D_m^-(a,b)=\sum_{n\ge0}\biggl\{&
\frac{x_n(-\log Q_b(x_n))^m}{Q_b(x_n)}
-\frac{(-2\log x_n)^m}{x_n}\notag\\[-2pt]
&-\frac{y_n(-\log Q_b(y_n))^m}{Q_b(y_n)}
+\frac{(-2\log y_n)^m}{y_n}\biggr\}.
\label{qspec:eq:Dminus}
\end{align}
The cancellation takes place separately at $x_n$ and $y_n$; no pairing of
two divergent sums is being used to define these differences.

\begin{theorem}[All-index, both-parity Stieltjes identity]
\label{qspec:thm:master}
For every $m\ge0$, $0<a<1$, and $b>0$,
\begin{align}
(-2)^mG_m^+(a)&=C_m(b)+B_m^+(a,b)-D_m^+(a,b),\label{qspec:eq:master-plus}\\
(-2)^mG_m^-(a)&=B_m^-(a,b)-D_m^-(a,b).\label{qspec:eq:master-minus}
\end{align}
In particular each individual Stieltjes constant is given by
\begin{equation}\label{qspec:eq:individual}
\gamma_m(a)=\frac{C_m+B_m^++B_m^--D_m^+-D_m^-}{2(-2)^m}.
\end{equation}
These identities may be differentiated any fixed number of times in $a$
and $b$ on compact subsets of \eqref{qspec:eq:domain}.
\end{theorem}
\begin{proof}
For $|\Rea\epsilon|<1/4$ put
\begin{align*}
D_+(\epsilon)&=\sum_{n\ge0}\left\{
Q_b(x_n)^{-1/2-\epsilon}-x_n^{-1-2\epsilon}
+Q_b(y_n)^{-1/2-\epsilon}-y_n^{-1-2\epsilon}\right\},\\
D_-(\epsilon)&=\sum_{n\ge0}\left\{
x_nQ_b(x_n)^{-1-\epsilon}-x_n^{-1-2\epsilon}
-y_nQ_b(y_n)^{-1-\epsilon}+y_n^{-1-2\epsilon}\right\}.
\end{align*}
Writing $Q_b(x)=x^2(1+b^2/x^2)$ shows that each local difference
is $O(x^{-3+2|\Rea\epsilon|})$, uniformly on compact sets of the
parameters and of $\epsilon$. Differentiation in $\epsilon$ introduces
at most fixed powers of $\log x$. The resulting series converge normally;
the same reasoning applies after finitely many $a,b$ derivatives.
Their $m$th derivatives at zero are exactly
\eqref{qspec:eq:Dplus} and~\eqref{qspec:eq:Dminus}.

In the initially convergent region, and hence as local meromorphic germs,
\begin{align}
E(\tfrac12+\epsilon)&=
\zeta(1+2\epsilon,a)+\zeta(1+2\epsilon,1-a)+D_+(\epsilon),\label{qspec:eq:even-subtraction}\\
Y(1+\epsilon)&=
\zeta(1+2\epsilon,a)-\zeta(1+2\epsilon,1-a)+D_-(\epsilon).
\label{qspec:eq:odd-subtraction}
\end{align}
Each individual Hurwitz term has residue $1/(2\epsilon)$, so the even
sum has residue $1/\epsilon$ and the odd difference has none. Equation
\eqref{qspec:eq:stieltjes-convention} gives the $m$th regular derivative
as $(-2)^mG_m^\pm$. Substitute \eqref{qspec:eq:germ-identities}, remove
the one explicit pole, and compare derivatives. This proves both
identities and their differentiated versions. Adding them proves
\eqref{qspec:eq:individual}.
\end{proof}

\subsection{All coefficients of the pole term}
The zero-mode coefficients have a particularly short finite description.
Set
\begin{equation}\label{qspec:eq:ell}
\ell_1=2\log(2/b),\qquad
\ell_r=(-1)^r(r-1)!(2-2^r)\zeta(r)\quad(r\ge2).
\end{equation}
\begin{proposition}\label{qspec:prop:pole-bell}
For every $m\ge0$,
\begin{equation}\label{qspec:eq:Cbell}
C_m(b)=\frac{\Bcal_{m+1}(\ell_1,\ldots,\ell_{m+1})}{m+1}.
\end{equation}
Writing $L=\log(2/b)$, the first three are
\begin{equation}\label{qspec:eq:Cexamples}
C_0=2L,\qquad C_1=2L^2-\zeta(2),\qquad
C_2=\frac83L^3-\frac{2\pi^2}{3}L+4\zeta(3).
\end{equation}
\end{proposition}
\begin{proof}
The function $\epsilon C(\epsilon;b)$ is analytic and equals one at zero.
Its logarithm is
\[
\log\Gamma(1+\epsilon)-\log\Gamma(1/2+\epsilon)
+\tfrac12\log\pi-2\epsilon\log b.
\]
The classical polygamma values at $1$ and $1/2$ \cite{gamma} give precisely
\eqref{qspec:eq:ell} as its successive derivatives at zero. Expand its
exponential by \eqref{qspec:eq:bell}, divide by $\epsilon$, remove the
pole, and multiply the coefficient of $\epsilon^m$ by $m!$. The resulting
factor is $m!/(m+1)!=1/(m+1)$, proving the proposition.
\end{proof}

\subsection{Explicit Fourier--Bessel spectral jets}
For $n\ge1$ set
\begin{align}
v_1(n)&=\gamma+\log(4\pi n/b),&
v_r&=(-1)^{r+1}(r-1)!(2^r-1)\zeta(r)\quad(r\ge2),\label{qspec:eq:vjets}\\
w_1(n)&=\gamma+\log(\pi n/b),&
w_r&=(-1)^{r+1}(r-1)!\zeta(r)\quad(r\ge2).\label{qspec:eq:wjets}
\end{align}
Define $K^{\langle r\rangle}_\nu(u)=\partial_\nu^rK_\nu(u)$.
Differentiating the holomorphic factors in
\eqref{qspec:eq:Bplusgerm}--\eqref{qspec:eq:Bminusgerm} gives
\begin{align}
B_m^+={}&4\sum_{n\ge1}\cos(2\pi na)
 \sum_{r=0}^{\lfloor m/2\rfloor}\binom{m}{2r}
 \Bcal_{m-2r}(v(n))K^{\langle2r\rangle}_0(2\pi bn),\label{qspec:eq:Bplusjets}\\
B_m^-={}&4\pi\sqrt b\sum_{n\ge1}\sqrt n\sin(2\pi na)
 \sum_{r=0}^m\binom mr\Bcal_{m-r}(w(n))
 K^{\langle r\rangle}_{-1/2}(2\pi bn).
\label{qspec:eq:Bminusjets}
\end{align}
Here $\Bcal_j(v(n))$ means that only the first $j$ entries are used.
Equation~\eqref{qspec:eq:K-cosh} shows that $K_\nu=K_{-\nu}$, so all
odd order derivatives at zero vanish. No corresponding deletion is valid
at $\nu=-1/2$. This explains the different inner sums.

For example,
\[
B_0^+=4\sum_{n\ge1}\cos(2\pi na)K_0(2\pi bn),\qquad
B_1^+=4\sum_{n\ge1}\cos(2\pi na)
\bigl(\gamma+\log(4\pi n/b)\bigr)K_0(2\pi bn).
\]
Thus \eqref{qspec:eq:master-plus} at $m=1$ is the explicit identity
\begin{align}
-2G_1^+(a)={}&2\log^2(2/b)-\zeta(2)-D_1^+(a,b)\notag\\
&+4\sum_{n\ge1}\cos(2\pi na)
\bigl(\gamma+\log(4\pi n/b)\bigr)K_0(2\pi bn).
\label{qspec:eq:first-stieltjes-example}
\end{align}
It involves no Bessel order derivative. Higher $m$ introduce them only
at the even orders displayed in \eqref{qspec:eq:Bplusjets}.

\subsection{Damped polylogarithm integral formulas}
Put
\begin{equation}\label{qspec:eq:zt}
z_t^\pm=\exp(-2\pi b\cosh t\pm2\pi ia),\qquad
\mathcal C_\sigma(t)=\frac{\Li_\sigma(z_t^+)+\Li_\sigma(z_t^-)}2,
\quad
\mathcal S_\sigma(t)=\frac{\Li_\sigma(z_t^+)-\Li_\sigma(z_t^-)}{2i}.
\end{equation}
For real $\sigma$, these are respectively the real and imaginary parts of
$\Li_\sigma(z_t^+)$. For complex $\sigma$, the symmetric formulas in
\eqref{qspec:eq:zt}, not a real-part operation, preserve holomorphy.

\begin{theorem}[Polylogarithm order integrals for both reflected channels]
\label{qspec:thm:polylog-integrals}
The entire germs $B_\pm$ satisfy
\begin{align}
B_+(\epsilon;a,b)&=
\frac{4\pi^{1/2+\epsilon}b^{-\epsilon}}{\Gamma(1/2+\epsilon)}
\int_0^\infty\cosh(\epsilon t)\mathcal C_{-\epsilon}(t)\dd t,
\label{qspec:eq:polylogplus}\\
B_-(\epsilon;a,b)&=
\frac{4\pi^{1+\epsilon}b^{1/2-\epsilon}}{\Gamma(1+\epsilon)}
\int_0^\infty\cosh((\epsilon-\tfrac12)t)
\mathcal S_{-\epsilon-1/2}(t)\dd t.
\label{qspec:eq:polylogminus}
\end{align}
Every spectral derivative in Theorem~\ref{qspec:thm:master} can therefore
be represented by differentiating an ordinary absolutely convergent
polylogarithm integral.
\end{theorem}
\begin{proof}
Insert \eqref{qspec:eq:K-cosh} in the Bessel series. The polylogarithm
series \cite{polylog} gives
\[
\mathcal C_{-\epsilon}(t)=\sum_{n\ge1}n^\epsilon
 e^{-2\pi bn\cosh t}\cos(2\pi na),
\]
with the analogous sine formula containing $n^{\epsilon+1/2}$.
For $\epsilon$ in a compact set, the terms and their derivatives have
majorants of the form
\[
C n^A(1+\log n)^M(1+t)^M e^{At}e^{-2\pi b_0n\cosh t},
\qquad b\ge b_0>0.
\]
Summing in $n$ and integrating in $t$ gives a finite quantity. The decay
in $t$ is stronger than any exponential. Fubini's theorem and holomorphic
parameter differentiation therefore justify the interchange and all
claimed spectral derivatives.
\end{proof}

\subsection{Zeroth-order reflection and an exact mass-transfer law}
Since $K_{1/2}(u)=\sqrt{\pi/(2u)}e^{-u}$,
\begin{equation}\label{qspec:eq:odd-zero}
B_0^-(a,b)=2\pi\sum_{n\ge1}e^{-2\pi bn}\sin(2\pi na)
=\frac{\pi\sin(2\pi a)}{\cosh(2\pi b)-\cos(2\pi a)}.
\end{equation}
The odd master identity recovers the classical digamma reflection law:
\begin{equation}\label{qspec:eq:zero-reflection}
\pi\cot(\pi a)=
\frac{\pi\sin(2\pi a)}{\cosh(2\pi b)-\cos(2\pi a)}-D_0^-(a,b).
\end{equation}
At $a=1/2$ every odd term is exactly zero, a structural check rather
than a near-zero numerical test.

Because the left sides of the master identities are independent of $b$,
we also obtain, for arbitrary $b,c>0$,
\begin{align}
D_m^+(a,b)-D_m^+(a,c)&=C_m(b)-C_m(c)+B_m^+(a,b)-B_m^+(a,c),\label{qspec:eq:transferplus}\\
D_m^-(a,b)-D_m^-(a,c)&=B_m^-(a,b)-B_m^-(a,c).
\label{qspec:eq:transferminus}
\end{align}
These are exact identities between two deformations with the Stieltjes
constant eliminated. They generate harmonic/polylogarithm identities
without retaining an unevaluated Stieltjes symbol.
